\documentclass[a4paper,11pt]{book}
\usepackage{listings}
\usepackage[utf8]{inputenc}
\usepackage{titlesec}
\usepackage{fancyhdr}
\usepackage[spanish,es-tabla]{babel}
\usepackage[hidelinks]{hyperref}
\usepackage{xcolor}
\usepackage{pdfpages}
\usepackage{eurosym}
\usepackage{graphicx}
\usepackage{caption}
\usepackage{subcaption}

% Información reutilizable
\newcommand{\asunto}{Trabajo de Fin de Grado}
\newcommand{\titulo}{Enfocador Astronomico Automático}
\newcommand{\tituloEng}{Automatic astronomical focuser}
\newcommand{\grado}{Grado en Ingeniería Informática}
\newcommand{\autor}{José Miguel López Pérez}
\newcommand{\email}{josmilope@gmail.com}
\newcommand{\tutor}{Dr. Sergio Alonso Burgos}
\newcommand{\escuela}{Escuela Técnica Superior de Ingenierías Informática y de Telecomunicación}
\newcommand{\universidad}{Universidad de Granada}
\newcommand{\ciudad}{Granada}
\newcommand{\vers}{Versión 1.0}

% Información archivo
\hypersetup{
	pdfauthor = {\autor\ (\email)},
	pdftitle = {\titulo},
	pdfsubject = {\asunto},
	pdfkeywords = {atronomía, hardware libre, procesamiento imágenes, INDI, arduino,  software libre},
	pdfcreator = {LaTeX con el paquete texlive},
	pdfproducer = {pdflatex}
}

% Estilo de cabeceras
\pagestyle{fancy}
\fancyhf{}
\fancyhead[LO]{\leftmark}
\fancyhead[RE]{\rightmark}
\fancyhead[RO,LE]{\textbf{\thepage}}
\setlength{\headheight}{1.5\headheight}

% Redefinición de comandos
\renewcommand{\lstlistingname}{Fragmento de código}
\renewcommand{\lstlistlistingname}{Índice de fragmentos de código}
\renewcommand{\chaptermark}[1]{\markboth{\textbf{#1}}{}}
\renewcommand{\sectionmark}[1]{\markright{\textbf{\thesection. #1}}}

% Definición de colores
\definecolor{gray97}{gray}{.97}
\definecolor{gray75}{gray}{.75}
\definecolor{gray45}{gray}{.45}
\definecolor{gray30}{gray}{.94}
\definecolor{lightgray}{rgb}{.9,.9,.9}
\definecolor{darkgray}{rgb}{.4,.4,.4}
\definecolor{purple}{rgb}{0.65, 0.12, 0.82}
\definecolor{background}{HTML}{EEEEEE}
\definecolor{delim}{RGB}{20,105,176}
\colorlet{punct}{red!60!black}
\colorlet{numb}{magenta!60!black}

% Listados
\lstset{
	aboveskip=0.5cm,
	backgroundcolor=\color{gray97},
	basicstyle=\scriptsize\ttfamily,
	breaklines=true,
	commentstyle=\color{gray45},
	frame=Ltb,
	framerule=0.5pt,
	framesep=0pt,
	framexbottommargin=3pt,
	framexleftmargin=0.1cm,
	framextopmargin=3pt,
	keywordstyle=\bfseries,
	numberfirstline = false,
	numbers=left,
	numbersep=6pt,
	numberstyle=\tiny,
	rulesep=.4pt,
	rulesepcolor=\color{black},
	showstringspaces = false,
	stringstyle=\ttfamily,
	literate={á}{{\'a}}1
	         {é}{{\'e}}1
	         {í}{{\'i}}1
	         {ó}{{\'o}}1
	         {ú}{{\'u}}1
	         {ñ}{{\~n}}1
}

% Minimizar fragmentado de listados
\lstnewenvironment{listing}[1][]
	{\lstset{#1}\pagebreak[0]}{\pagebreak[0]}

% Listado definido para JavaScript
% http://tex.stackexchange.com/questions/89574/language-option-supported-in-listings/89576#89576
\lstdefinelanguage{javascript}{
	backgroundcolor=\color{background},
	basicstyle=\footnotesize,
	breaklines=true,
	captionpos=b,
	comment=[l]{//},
	commentstyle=\color{purple}\ttfamily,
	frame=lines,
	identifierstyle=\color{black},
	keywordstyle=\color{blue}\bfseries,
	morecomment=[s]{/*}{*/},
	morestring=[b]',
	morestring=[b]",
	ndkeywordstyle=\color{darkgray}\bfseries,
	numbers=left,
	numbersep=8pt,
	numberstyle=\scriptsize,
	sensitive=false,
	showstringspaces=false,
	stepnumber=1,
	stringstyle=\color{red}\ttfamily,
	keywords={
		break,
		case,
		catch,
		catch,
		do,
		else,
		false,
		function,
		if,
		in,
		new,
		null,
		return,
		switch,
		true,
		typeof,
		var,
		while},
	ndkeywords={
		boolean,
		class,
		export,
		implements,
		import,
		this,
		throw}
}

% Listado definido para JSON
% http://tex.stackexchange.com/questions/83085/how-to-improve-listings-display-of-json-files/83100#83100
\lstdefinelanguage{json}{
	backgroundcolor=\color{background},
	basicstyle=\footnotesize,
	breaklines=true,
	captionpos=b,
	frame=lines,
	numbers=left,
	numbersep=8pt,
	numberstyle=\scriptsize,
	showstringspaces=false,
	stepnumber=1,
	literate=
		*{:}{{{\color{punct}{:}}}}{1}
		{,}{{{\color{punct}{,}}}}{1}
	    {\{}{{{\color{delim}{\{}}}}{1}
	    {\}}{{{\color{delim}{\}}}}}{1}
	    {[}{{{\color{delim}{[}}}}{1}
	    {]}{{{\color{delim}{]}}}}{1}
	    {ñ}{{\~{n}}}{1}
}

% Para que las páginas en blanco no tengan cabecera
\makeatletter
\def\clearpage{%
  \ifvmode
    \ifnum \@dbltopnum =\m@ne
      \ifdim \pagetotal <\topskip
        \hbox{}
      \fi
    \fi
  \fi
  \newpage
  \thispagestyle{empty}
  \write\m@ne{}
  \vbox{}
  \penalty -\@Mi
}
\makeatother

\begin{document}
\begin{titlepage}
Hola travis
\newlength{\centeroffset}
\setlength{\centeroffset}{-0.5\oddsidemargin}
\addtolength{\centeroffset}{0.5\evensidemargin}

\noindent\hspace*{\centeroffset}\begin{minipage}{\textwidth}

\centering
\includegraphics[width=0.9\textwidth]{../images/logo_ugr.png}\\[1.4cm]

\textsc{\Large\asunto\\[0.2cm]}
\textsc{\grado}\\[1cm]

{\Huge\bfseries \titulo\\}
\noindent\rule[-1ex]{\textwidth}{3pt}\\[3.5ex]
\end{minipage}

\vspace{2cm}
\noindent\hspace*{\centeroffset}\begin{minipage}{\textwidth}
\centering

\textbf{Autor}\\ {\autor}\\[2.5ex]
\textbf{Tutor}\\ {\tutor}\\[2cm]
\includegraphics[width=0.3\textwidth]{../images/logo_etsiit.png}\\[0.1cm]
\textsc{\escuela}\\
\textsc{---}\\
\ciudad, \today\\
\includegraphics[width=0.3\textwidth]{../images/CC-SA-logo.png}
\end{minipage}
\end{titlepage}

\frontmatter
\begin{center}
{\LARGE\bfseries\titulo}\\
\end{center}
\begin{center}
\autor\
\end{center}

\begin{figure}[h]
\centering
\includegraphics[width=0.7\linewidth]{../images/simple}
\caption[Ardufocuser]{Ardufocuser}
\label{fig:logo}
\end{figure}

\newpage




\begingroup
\let\cleardoublepage\clearpage
  \tableofcontents
  \listoffigures
  \listoftables
  \lstlistoflistings
\endgroup

\newpage
\thispagestyle{empty}
\
\mainmatter
\chapter{Resumen}

\section{Breve resumen y palabras clave}
\noindent{\textbf{Palabras clave}: \textit{arduino}, \textit{hardware libre}, \textit{astronomía}, \textit{procesamiento imagen}, \textit{}, \textit{INDI}, \textit{software libre}.\\

\bigskip
El objetivo principal de este proyecto es diseñar e implementar un sistema completo de enfoque automático, que se pueda acoplable directamente
 a un telescopio comercial y consiga de forma óptima realizar una configuración de las lentes para conseguir el mejor enfoque posible de las imágenes estelares.

Para ello debemos hacer un recorrido tecnológico que pase por el diseño y programación a bajo nivel de un dispositivo hardware que interactúe físicamente con el telescopio,
así como ocuparnos de la comunicación con el servidor al que se conectan los clientes distribuidos, las métricas y heurísticas para evaluar la calidad de las imagenes que se toman,
así como los distintos algoritmos para encontrar los máximos deseados. Topándonos de lleno con un amplio abanico  aspectos relacionados con el diseño, la computación, desarrollo software
y las comunicaciones.

\bigskip
Dado la gran historia y recorrido de la astronomía, podemos apreciar una gran evolución en las herramientas astronómicas que han ido apareciendo,
 sirviéndose siempre de los últimos avances en los demás campos de ciencias como puede ser la óptica, las matemáticas, la mecánica, la electrónica y la informática.

\bigskip
Hoy día ya contamos con numerosas plataformas de control y automatización de dispositivos astronómicos, aunque en muchos casos se relacionan con grandes compañías propietarias
de gran cantidad de derechos intelectuales, con el inconveniente de no ofrecen detalles técnicos, información bien detallada y libertad para conocer el funcionamiento y
poder modificarlo.

\bigskip
El reto al que nos enfrentamos en este proyecto es acercar  hardware y el software libre a la astronomía creando herramientas avanzadas con un bajo presupuesto
y alcance libre de todas sus interioridades, detalles técnicos así como la libertad para mejorar y evolucionar el producto por todo aquel interesado.

Permitiendo que este proyecto además de ser un producto perfectamente funcional, sea una semilla para multitud de nuevos proyectos en el contexto de la astronomía.


\newpage
\begin{center}
{\LARGE\bfseries\tituloEng}\\
\end{center}
\begin{center}
\autor\
\end{center}

\section{Extended abstract and key words}

\noindent{\textbf{Palabras clave}: \textit{arduino}, \textit{hardware libre}, \textit{astronomía}, \textit{procesamiento imagen}, \textit{}, \textit{INDI}, \textit{software libre}.\\


\bigskip
The main objective of this project is to design and implement a complete autofocus system that can be coupled directly to a commercial telescope and get optimally perform a configuration of lenses, to observe stellar images.



\backmatter

\newpage
\begin{thebibliography}{99}
	\addcontentsline{toc}{chapter}{Bibliografía}


\bibitem{INDILIB} ``INDI'', Fecha de último acceso: 1/9/2015.\\ \href{http://indilib.org/}{http://indilib.org/}

\end{thebibliography}

\newpage \
\thispagestyle{empty}
\end{document}
